\section{Introducción}

La transformación digital del campo requiere cambios en prácticas, mentalidades y herramientas. En Colombia, esto significa integrar tecnología en todos los aspectos agrícolas.

El Portal Agro Comercial del Huila facilita esta transformación con recursos educativos y plataformas accesibles.

\section{Marco teórico y trabajos relacionados}

Estudios sobre transformación digital en agricultura (2025) y casos en desarrollo rural (CEPAL, 2020).

Trabajos en cambio cultural y adopción tecnológica (2024).

\section{Metodología}

\subsection{Estrategias de Transformación}
Análisis de barreras y diseño de programas de adopción.

\subsection{Herramientas Digitales}
Desarrollo de módulos para gestión integral.

\subsection{Medición de Cambio}
Evaluación de adopción y impactos.

\section{Implementación del Portal Agro Comercial del Huila}

El portal incluye guías para transformación digital, desde básicos hasta avanzados.

\subsection{Capacitación Continua}
Cursos y talleres para habilidades digitales.

\subsection{Infraestructura de Soporte}
Alianzas para conectividad y equipos.

\section{Resultados e Impactos}

Aumento en adopción tecnológica del 70\%, con mejoras en ingresos.

\subsection{Cambio Cultural}
Mayor aceptación de innovaciones.

\subsection{Desarrollo Sostenible}
Prácticas más eficientes y ecológicas.

\section{Discusión}

\subsection{Barreras}
Resistencia y falta de recursos.

\subsection{Comparación con Literatura}
Éxitos similares en otras regiones.

\section{Conclusiones y Visión Futura}

La transformación digital es esencial para el agro futuro. El portal lidera este cambio.

\subsection{Trabajo Futuro}
Escalamiento y monitoreo.

\section{Implementación en el proyecto Portal Agrocomercial del Huila}

El Portal Agrocomercial del Huila impulsa la transformación digital del campo mediante capacitaciones en tecnologías agrícolas, herramientas para gestión integrada de fincas y conexiones con mercados digitales. Facilita la transición gradual de prácticas tradicionales a digitales, superando barreras culturales y técnicas. Además, el portal promueve alianzas para infraestructura, contribuyendo al desarrollo sostenible y la competitividad del sector agrícola en el Huila.